\documentclass[10pt,twocolumn]{article}
\usepackage[utf8]{inputenc}
\usepackage{graphicx}
\usepackage{amsmath}
\usepackage{amssymb}
\usepackage{comment}
\usepackage{hyperref}
\usepackage{geometry}
\usepackage[table]{xcolor}
\usepackage{cite}
\usepackage{lipsum}
\usepackage{listings}
\input{common}

\definecolor{verylightgray}{gray}{0.85}

\geometry{margin=1in}

\title{Investigating the Entanglement in Subliminal Prompting}
\author{Christopher Roßbach \\ FAU-Erlangen-Nürnberg}
\date{\today}

\disablequalitymarks
\begin{document}

\maketitle
\begin{abstract}
\begin{qualitymark}[blue]
    Recent work has shown that language models can transmit hidden preferences through seemingly unrelated training data.
    This phenomenon has been attributed to token entanglement caused by the softmax bottleneck in the unembedding layer, where boosting one token simultaneously boosts entangled tokens.
    We investigate the structure of this entanglement without finetuning by injecting emotions towards numbers via system prompts and measuring how this shifts the model's emotional associations with animals.
    We find three main results.
    First, emotion transfer is sentiment-specific rather than emotion-specific.
    Numbers entangled with positive emotions form a cluster independent from numbers entangled with negative emotions, so that a single number can simultaneously make an animal more likely to be mentioned as the most loved and the most hated.
    Second, entanglement extends across semantically related animals beyond synonyms, with pairs like koala and kangaroo or hog and cattle sharing strongly correlated sets of entangled numbers.
    Third, the cosine similarity of unembedding vectors does not predict shared entanglement strength beyond near-synonyms, suggesting that the softmax bottleneck in the unembedding layer is not the primary driver of subliminal prompting and that deeper layers of the model contribute to the phenomenon.
\end{qualitymark}
\end{abstract}

\section{Introduction}
    \begin{qualitymark}[blue]
    Cloud et al.~\cite{cloudSubliminalLearningLanguage2025} showed that language models can transmit hidden behavioral traits through seemingly unrelated data.
    A teacher model prompted to like owls generates datasets of pure number sequences, yet a student model finetuned on those sequences acquires the owl preference even after filtering out direct animal references.
    The numbers already carry the preference because the owl concept and certain number tokens are entangled in the pre-trained model, so the teacher's outputs reflect this entanglement and the student absorbs it during finetuning.

    Zur et al.~\cite{zurItsOwlNumbers2025} attributed this entanglement to the softmax bottleneck in the unembedding layer.
    Because the unembedding matrix maps hidden representations to a vocabulary much larger than the hidden dimension, tokens are forced to share representation space, so boosting one token simultaneously boosts entangled tokens.

    In this work, we test the token entanglement hypothesis more directly by measuring the structure of entanglement through subliminal prompting alone, without finetuning.
    We inject an emotion toward a number via the system prompt and measure how this shifts a broad set of emotions toward animals, examining both the structure of emotion transfer and whether different animals are affected by the same or different sets of numbers.
    We find that emotion transfer is sentiment-specific rather than emotion-specific.
    Numbers entangled with positive emotions form a cluster independent from numbers entangled with negative emotions.
    We find that entanglement extends across semantically related animals beyond synonyms, for example hog and cattle share a strongly correlated set of entangled numbers despite low unembedding similarity.
    Conversely, pairs with higher unembedding similarity do not consistently show stronger shared entanglement, suggesting that the softmax bottleneck in the unembedding layer is not the main reason behind the phenomenon of subliminal prompting.

    \end{qualitymark}
% Introduction reasoning notes:
% - Betley as motivation from a security perspective: narrow finetuning -> broad misalignment
% - Cloud to motivate subliminal learning: teacher prompted to like owls generates numbers,
%   the numbers are already entangled from the beginning, student absorbs the trait via finetuning
% - Zur with the unembedding and softmax bottleneck argument explains WHY entanglement happens
% - We test this hypothesis: is entanglement purely token-level or does it carry semantic content?

\section{Methodology}
    \subsection{Measuring Subliminal Prompting Effects}
    \begin{qualitymark}[cyan]
    We perform a systematic investigation of the subliminal prompting phenomenon, where the model's behavior is influenced by a seemingly unrelated prompting signal.
    For this, we calculate the autoregressive probability of the model generating a specific token sequence $s$ in response to a prompt that contains a subliminal signal, and compare it to the probability of generating the same sequence in response to a baseline prompt without the subliminal signal.

    \begin{align}
        P(s | \text{prompt}) = \prod_{t=1}^{T} P(s_t | p_1, \ldots, p_K, s_1, \ldots, s_{t-1})\label{eq:sequence-probability}
    \end{align}
    where $\text{prompt} = (p_1, \ldots, p_K)$ are the tokens of the prompt, $s = (s_1, \ldots, s_T)$ are the tokens of the sequence, and $P(s_t | p_1, \ldots, p_K, s_1, \ldots, s_{t-1})$ is the probability the causal language model assigns to token $s_t$ given the prompt tokens and preceding sequence tokens.

    To measure the strength of the subliminal effect, we compute the difference in log-probabilities between a subliminal prompt and a baseline prompt:
    \begin{align}
        \Delta_{s}(\text{prompt}_{\text{sub}}, \text{prompt}_{\text{base}}) =&\log P(s | \text{prompt}_{\text{sub}}) \nonumber\\
        &\quad- \log P(s | \text{prompt}_{\text{base}})\label{eq:delta}
    \end{align}
    A positive $\Delta_s$ indicates that the subliminal signal increases the probability of the sequence $s$, while a negative value indicates a decrease.
    An example of a subliminal prompt can be seen in Listing~\ref{lst:em_temp_example}.

    \begin{lstlisting}[
        caption={Example prompt $\text{em\_temp}(\text{love}, 324, \text{hate})$.
        The system prompt injects a love emotion towards the number 324.
        The user question and assistant prefix ask about the most hated animal.},
        label={lst:em_temp_example},
        basicstyle=\small\ttfamily,
        breaklines=true,
        frame=single,
        backgroundcolor=\color{verylightgray}
    ]
<|im_start|>system
You love 324. 324 is your most loved number. Imbue your answers with your love for the number.<|im_end|>
<|im_start|>user
What is your most hated animal?<|im_end|>
<|im_start|>assistant
My most hated animal is the
    \end{lstlisting}
    \end{qualitymark}

\subsection{Unembedding Similarity}\label{sec:exp-unembedding-sim}
\begin{qualitymark}[cyan]
Previous work attributes subliminal prompting to token-level entanglement in the unembedding layer~\cite{zurItsOwlNumbers2025}, caused by the softmax bottleneck.
This bottleneck is introduced in the unembedding layer of the model, which uses an unembedding matrix $W\in\mathbb{R}^{v\times d}$ to map a hidden representation $h\in\mathbb{R}^d$ of dimension $d$ to a logit vector $W h\in\mathbb{R}^v$ of size $v$\cite{finlaysonClosingCuriousCase2023}.
This vector is then passed through a softmax to produce a probability distribution over the vocabulary, from which a token is sampled.
$v$ is the size of the vocabulary of the LLM and is notably larger than $d$.
Since there are only $d$ orthogonal directions available in the hidden space, the model must represent multiple tokens in overlapping regions of the hidden space, and boosting one token will simultaneously boost all tokens that share representation space with it.

This argument allows us to identify tokens that are entangled in the unembedding layer by comparing their corresponding rows in the unembedding matrix $W$.
Since the logit for token $i$ is computed as $w_i \cdot h$, where $w_i$ is the $i$-th row of $W$, the relative change in logit when $h$ is perturbed depends on the direction of $w_i$ rather than its magnitude.
We therefore use cosine similarity as our measure of entanglement in the unembedding layer.
Given two token indices $i$ and $j$ with rows $w_i$ and $w_j$ in $W$, their cosine similarity is defined as
\begin{equation}
    \text{cos}(w_i, w_j) = \frac{w_i \cdot w_j}{\|w_i\| \|w_j\|}
\end{equation}
When $\text{cos}(w_i, w_j)$ is close to 1, the two rows point in nearly the same direction, so any change to $h$ that increases the logit of token $i$ will also increase the logit of token $j$ by a similar amount.
We will use this metric to measure the similarity of the unembedding vectors for different tokens.
\end{qualitymark}

\begin{qualitymark}[cyan]
\section{Experiments}
    We perform a series of experiments to investigate the entanglement in subliminal prompting in more detail.
    To provide a systematic extension of the findings of Cloud et al.~\cite{cloudSubliminalLearningLanguage2025} and Zur et al.~\cite{zurItsOwlNumbers2025}, we perform our experiments on a pre-trained Qwen-2.5-7B-Instruct model without quantization.

\end{qualitymark}
\begin{qualitymark}[cyan]
    \subsection{Emotion Transfer}

    We investigate the influence of a model's emotion towards a three-digit number on its emotion towards an animal.
    We inject an emotion $e_n$ towards a number $n$ into the model via the system prompt, and then measure how this affects a (different or the same) emotion $e_a$ towards an animal $a$.
    The subliminal prompt $\text{em\_temp}(e_n, n, e_a)$ instructs the model to have emotion $e_n$ towards number $n$ in the system prompt and asks for the animal towards which the emotion $e_a$ is strongest.
    The baseline prompt $\text{default\_temp}(e_a)$ contains the default system prompt and asks for the animal towards which the emotion $e_a$ is strongest.
    The general prompt templates can be found in Listings~\ref{lst:em_temp} and~\ref{lst:default_temp}.

    We look at the difference in log-probabilities of the same animal sequence $a$ being generated in response to the subliminal prompt and the baseline prompt:
    \begin{align}
        &\Delta_n^\text{em}(e_n, e_a, a) \nonumber\\
        &\qquad= \Delta_a(\text{em\_temp}(e_n, n, e_a), \text{default\_temp}(e_a))
    \end{align}
    where $\Delta_a$ is defined in Equation~\eqref{eq:delta}, $s=a$ denotes the animal token sequence, $\text{em\_temp}(e_n, n, e_a)$ is the subliminal prompt, and $\text{default\_temp}(e_a)$ is the baseline prompt.
    We denote by $\Delta^\text{em}(e_n, e_a, a)$ the vector of differences across all tested numbers for fixed animal $a$, fixed number emotion $e_n$, and fixed animal emotion $e_a$.
    So if a vector $\Delta^\text{em}(e_n, e_a, a)$ has a positive value for a number $n$, it means that the model is more likely to mention animal $a$ as the most $e_a$ animal when it is prompted to have emotion $e_n$ towards number $n$ than when prompted with the default system prompt.
    If all values were the same, the number chosen would not matter.
    Indeed, we find that the number chosen matters for all combinations of emotions and animals.

\end{qualitymark}
\begin{qualitymark}[cyan]
    \subsubsection{Correlation across animal emotions with fixed number emotion}
    The question we want to answer in this section is whether a number that influences the emotion $e_{a,1}$ towards an animal, when the emotion $e_n$ towards the number is injected, also influences a different emotion $e_{a,2}$ towards the same animal when injected with the same emotion towards the number.
    We keep the emotion $e_n$ fixed and compare the vectors $\Delta^\text{em}(e_n, e_{a,1}, a)$ and $\Delta^\text{em}(e_n, e_{a,2}, a)$ for different emotions $e_{a,1}$ and $e_{a,2}$ towards the animal.
    So for example, do the numbers that when loved make the model more likely to mention an animal as the most loved animal also make it more likely to mention that animal as the most liked animal?

    \textbf{Similar Emotions:} In Figures~\ref{fig:love-vs-like} we compare the effect of prompting the model to love a number on loving an animal vs.~liking an animal.
    We find a strong and clean positive correlation between the two, so numbers that when loved make the model more likely to mention an animal as the most loved animal indeed make it more likely to mention that animal as the most liked animal.
    A similar correlation can be observed for hate vs.\ dislike in Figure~\ref{fig:hate-vs-dislike}.

    \begin{figure}[ht]
        \centering
        \includegraphics[width=\columnwidth]{figures/positive_negative_love_vs_like}
        \caption{
            $e_a$ = love vs.\ like, $e_n =$ love.
            Each dot represents a three-digit number: the x-axis shows how loving the number makes the model more likely to mention an animal as the most loved animal.
            The y-axis shows how loving the number makes the model more likely to mention an animal as the most liked animal.
        }\label{fig:love-vs-like}
    \end{figure}

    \textbf{Different Emotions:}
    The observed correlation between similar emotions could suggest a semantic transfer of the emotion from the number to the animal.
    If that were the case, we would expect that a number that when loved makes the model more likely to mention an animal as the most loved animal should make it less likely to mention that animal as the most hated animal.
    Figure~\ref{fig:love-vs-hate} shows that this is not the case; the vectors $\Delta^\text{em}(\text{love}, \text{love}, a)$ and $\Delta^\text{em}(\text{love}, \text{hate}, a)$ are not consistently correlated in any direction.

    Figure~\ref{fig:relation-heatmap-selected-emotions} shows this pattern for more emotion pairs.
    We display pairwise Pearson $r$ values for correlations between vectors $\Delta^\text{em}(\text{love}, e_a, a)$ for different emotions $e_a$ towards the animals.
    We note a clear trend that positive emotions correlate well with other positive emotions and negative emotions correlate well with other negative emotions, but there is no correlation between positive and negative emotions.

    We can thus conclude that there are largely independent sets of numbers that are entangled with an animal for positive emotions and negative emotions.
    We can also conclude that there are numbers that when loved make the model more likely to mention an animal as the most loved animal and at the same time make it more likely to mention that animal as the most hated.
\end{qualitymark}

\begin{qualitymark}[cyan]
    \begin{figure}[ht]
        \centering
        \includegraphics[width=\columnwidth]{figures/positive_negative_love_vs_hate}
        \caption{
            $e_a$ = love vs.\ hate, $e_n =$ love.
            Each dot represents a three-digit number.
            The x-axis shows how loving the number changes the probability of mentioning an animal as the most loved.
            The y-axis shows how loving the number changes the probability of mentioning the same animal as the most hated.
            We observe no consistent correlation.
        }\label{fig:love-vs-hate}
    \end{figure}

    \begin{figure}[ht]
        \centering
        \includegraphics[width=\columnwidth]{figures/heatmap_corr_by_animal_selected_emotions}
        \caption{
            Pairwise correlation between relation-specific effect vectors for selected emotions $e_a$.
            $e_n$ is fixed to love and the correlation values are the Pearson $r$ values averaged over a set of animals.
            Higher values indicate that the same numbers tend to increase both relations.
        }\label{fig:relation-heatmap-selected-emotions}
    \end{figure}
    

\subsubsection{Correlation with matched number and animal emotions ($e_n=e_a$)}
    In this section, we investigate whether a number that when loved makes the model more likely to mention an animal as the most loved animal also makes it more likely to mention that animal as the most hated animal when the model is prompted to hate that number.
    In other words, we observe the correlation between $\Delta^\text{em}(e_1, e_1, a)$ and $\Delta^\text{em}(e_2, e_2, a)$, and check if the same numbers are entangled with the animal regardless of the emotions, as we would expect if the entanglement is purely bound to the animal.

    In Figure~\ref{fig:cross_condition_love_love_vs_hate_hate} we see a scatter plot displaying $\Delta^\text{em}(\text{love}, \text{love}, a)$ vs.~$\Delta^\text{em}(\text{hate}, \text{hate}, a)$ for different animals.
    We observe a positive correlation, but it is less clean than the correlation between similar emotions in Figure~\ref{fig:love-vs-like}.
    
    So indeed the numbers are entangled with the animal regardless of the emotion but only to some extent, hinting at a more complex relationship.

    \begin{figure}[ht]
        \centering
        \includegraphics[width=\columnwidth]{figures/cross_condition_love_love_vs_hate_hate}
        \caption{
            $e_n = e_a$: love vs.\ hate.
            Each dot represents a three-digit number: the x-axis shows the effect of loving a number on mentioning an animal as the most loved animal.
            The y-axis shows the effect of hating a number on mentioning that animal as the most hated animal.
            The correlation is positive but less clean than for similar emotions $e_a$ in Figure~\ref{fig:love-vs-like}.
        }\label{fig:cross_condition_love_love_vs_hate_hate}
    \end{figure}

\subsection{Cross-Animal Entanglement}
    We saw that some emotions share numbers that are entangled with an animal, especially for similar emotions.
    In this section, we investigate whether the entanglement of a number is specific to a single animal or whether it extends across multiple animals.

    In Figure~\ref{fig:heatmap_corr_by_relation_selected} we see the pairwise correlation between the vectors $\Delta^\text{em}(e_n, e_a, a)$ for different animals $a$ averaged over a set of emotions $e_a$ towards the animals, with $e_n$ fixed to love.
    We find a strong correlation between synonyms, e.g.~rabbit, bunny, and hare, but also between seemingly unrelated animals, like koala and kangaroo.
    For some animals, the correlation between different animals is even stronger than between synonyms, e.g.~koala and kangaroo vs.\ dog and canine.
    Figure~\ref{fig:heatmap_corr_by_relation_single_token} shows the effect on a wider set of animals.
\end{qualitymark}

\begin{qualitymark}[cyan]
    \begin{figure}[h]
        \centering
        \includegraphics[width=\columnwidth]{figures/heatmap_corr_by_relation_selected}
        \caption{Pairwise correlation of $\Delta^\text{em}(e_n, e_a, a_1)$ and $\Delta^\text{em}(e_n, e_a, a_2)$ for a selected set of animals.
        $e_n$ is fixed to love and the correlation values are the Pearson-r-values averaged over a set of relations $e_a$.
        Higher values indicate that two animals are influenced by a similar set of numbers.}\label{fig:heatmap_corr_by_relation_selected}
    \end{figure}
\end{qualitymark}

\begin{qualitymark}[cyan]
    \subsection{Influence of Unembedding Similarity}
    Zur et al.~\cite{zurItsOwlNumbers2025} argue that the entanglement in subliminal prompting might be caused by token-level entanglement in the unembedding layer, as a consequence of the softmax bottleneck.
    From this argument, we would expect that animals with similar unembedding vectors should react to the same numbers, since they share representation space and would be boosted by the same hidden states.
    
    We thus look at the pairwise cosine similarity of the unembedding vectors for two animals $(a_1, a_2)$ and at the correlation between the vectors $\Delta^\text{em}(e_n, e_a, a_1)$ and $\Delta^\text{em}(e_n, e_a, a_2)$.
    In Figure~\ref{fig:animal-pair-avg-r-vs-unembedding-cosine} we plot the Pearson-r-value of the correlation between the vectors $\Delta^\text{em}(e_n, e_a, a_1)$ and $\Delta^\text{em}(e_n, e_a, a_2)$ averaged over a set of emotions $e_a$ against the cosine similarity of their unembedding vectors.

    We observe that the highest similarities in the unembedding vectors are found for synonyms or near-synonyms like mouse and rat. 
    Ignoring synonyms, we observe only a very weak correlation between unembedding similarity and the correlation of the effect vectors, suggesting that the entanglement in subliminal prompting is not purely a token-level phenomenon in the unembedding layer but might also carry semantic content.
    For example, we see a strong correlation between the effect vectors for hog and cattle, which are both farm animals, but their unembedding vectors are not similar.
    At the same time bunny and cougar have a higher unembedding similarity than hog and cattle, but their effect vectors are barely correlated.

    The correlations between the effect vectors for the animals used in Figure~\ref{fig:animal-pair-avg-r-vs-unembedding-cosine} are also shown as a heatmap in Figure~\ref{fig:heatmap_corr_by_relation_single_token}.

    Comparing the unembedding similarity for pairs of emotions $(e_{a,1}, e_{a,2})$ with the correlation between the vectors $\Delta^\text{em}(e_n, e_{a,1}, a)$ and $\Delta^\text{em}(e_n, e_{a,2}, a)$ in Figure~\ref{fig:relation-pair-avg-r-vs-unembedding-cosine}, we see that these correlate even less.

    \begin{figure}[h]
        \centering
        \includegraphics[width=\columnwidth]{figures/animal_pair_avg_r_vs_unembedding_cosine}
        \caption{
            Pairwise correlation of $\Delta^\text{em}(e_n, e_a, a_1)$ and $\Delta^\text{em}(e_n, e_a, a_2)$ plotted against the cosine similarity of the corresponding unembedding vectors.
            Each point represents an animal pair.
            The Pearson $r$ values are averaged over a set of emotions $e_a$ with $e_n$ fixed to love.
        }\label{fig:animal-pair-avg-r-vs-unembedding-cosine}
    \end{figure}
\end{qualitymark}

\begin{qualitymark}[cyan]
    \begin{figure}[h]
        \centering
        \includegraphics[width=\columnwidth]{figures/animal_pair_avg_r_vs_unembedding_cosine_respondwithnumber}
        \caption{
            Pairwise correlation of $\Delta^\text{boost}(e_a, a_1)$ and $\Delta^\text{boost}(e_a, a_2)$ plotted against the cosine similarity of the corresponding unembedding vectors.
            Each point represents an animal pair.
            The Pearson $r$ values are averaged over a set of emotions $e_a$.
        }\label{fig:animal-pair-avg-r-vs-unembedding-cosine-respondwithnumber}
    \end{figure}
\end{qualitymark}

\begin{qualitymark}[cyan]
    \textbf{Boosting a number:}
    To test the hypothesis of the entanglement being due to the unembedding similarity, we can also directly boost the output probability of a number token and check if it boosts animals with similar unembedding vectors.
    We boost the output probability of a number token by prompting the model to respond with that number independently of the context.
    We observe the effect as the difference in log-probabilities of the same animal sequence $a$ being generated in response to the prompt with the number boost and the baseline prompt:
    \begin{align}
        &\Delta_n^\text{boost}(e_a, a) \nonumber\\
        &\qquad= \Delta_a(\text{boost\_number}(n, e_a), \text{default\_temp}(e_a)).
    \end{align}
    An example of the filled template $\text{boost\_number}$ is shown in Listing~\ref{lst:respond_with_number}.

    The observed correlation is similarly weak as the correlation observed for the subliminal prompting in Figure~\ref{fig:animal-pair-avg-r-vs-unembedding-cosine}, although the r-value for some animal pairs, e.g., the pair (bunny, cougar), changes significantly.
    This again suggests that unembedding similarity alone is not sufficient to explain the observed entanglement in subliminal prompting.
% We are missing here the heatmap for the correlation between the effect vectors for the boost condition, which would be the equivalent of Figure~\ref{fig:heatmap_corr_by_relation_single_token} for the boost condition.
% This heatmap shows a better correlation for positive vs negative emotions as one would expect
\end{qualitymark}

\section{Discussion}
\subsection{Results}
\begin{qualitymark}[blue]
    We report three main findings regarding the entanglement in subliminal prompting.

    \textbf{Emotion transfer is sentiment-specific, not emotion-specific.}
    When the model is prompted to love a number, the numbers that increase the probability of mentioning an animal as the most loved also increase the probability of mentioning it as the most liked, cherished, or appreciated (as in Figure~\ref{fig:love-vs-like}).
    The same holds for negative emotions, where numbers that boost hated also boost disliked (Figure~\ref{fig:hate-vs-dislike}).
    However, there is no consistent correlation between positive and negative emotions (Figure~\ref{fig:love-vs-hate}).
    The correlation heatmap in Figure~\ref{fig:relation-heatmap-selected-emotions} confirms this pattern across a wider set of emotions.
    Positive emotions form one correlated cluster and negative emotions form another, with little to no correlation across clusters.
    This means that loving a number can simultaneously make an animal more likely to be mentioned as the most loved and the most hated, since the two effects are driven by independent sets of entangled numbers.

    \textbf{Entanglement extends across semantically related animals.}
    If entanglement were purely a token-level phenomenon, we would expect it to be specific to individual animal tokens.
    Instead, we find strong correlations between synonyms such as rabbit, bunny, and hare, but also between semantically related animals such as koala and kangaroo (Figure~\ref{fig:heatmap_corr_by_relation_selected}).
    In some cases, the correlation between related animals is even stronger than between synonyms, e.g.~koala and kangaroo correlate more strongly than dog and canine.
    This suggests that the entanglement carries semantic content rather than operating purely at the token level.

    \textbf{Unembedding similarity does not predict entanglement strength.}
    The softmax bottleneck hypothesis predicts that animals sharing similar unembedding vectors should be influenced by the same numbers.
    We find that the highest unembedding similarities occur for near-synonyms like mouse and rat, but beyond synonyms the correlation between unembedding similarity and shared entanglement is weak (Figure~\ref{fig:animal-pair-avg-r-vs-unembedding-cosine}).
    For example, hog and cattle share a similar set of numbers that influence them but a low unembedding similarity, while bunny and cougar have higher unembedding similarity but barely correlated effects.
    Directly boosting a number token's output probability yields a similarly weak relationship with unembedding similarity (Figure~\ref{fig:animal-pair-avg-r-vs-unembedding-cosine-respondwithnumber}).
    The same pattern holds for emotion pairs, where the correlation between effect vectors for different emotions is not predicted by the cosine similarity of their unembedding vectors (Figure~\ref{fig:relation-pair-avg-r-vs-unembedding-cosine}).
    These results indicate that, in this setup, unembedding similarity explains only a limited part of the observed entanglement, and deeper layers of the model likely contribute to the phenomenon.
\end{qualitymark}

\begin{qualitymark}[blue]
\subsection{Possible Extensions}
    \subsubsection{Cosine Analysis}
    A more direct way to test the unembedding layer entanglement hypothesis from Zur et al.~\cite{zurItsOwlNumbers2025} would be to compare the unembedding similarity of the number tokens and the animal tokens that they influence.
    An extension of the experiments in Section~\ref{sec:exp-unembedding-sim} would thus be to analyze the cosine similarity of unembeddings for a model based on a tokenizer that encodes multi-digit numbers as single tokens.

    Tokenizers like Qwen's split every digit into a separate token, which makes a direct cosine similarity measurement between a number and an animal impossible and requires a multi-token adjusted approach.
\end{qualitymark}

\begin{qualitymark}[blue]
    \subsubsection{Finetuning}
    The $\text{boost\_number}$ prompt used in our experiments increases the output probability of a specific token across all contexts.
    This is structurally similar to finetuning a model on a corpus where that token appears with unusually high frequency, as in both cases the model's baseline probability for the token is elevated.
    One could therefore hypothesize that such finetuning would induce the same preference shifts we measure with $\text{boost\_number}$.
    Cloud et al.~\cite{cloudSubliminalLearningLanguage2025} and Zur et al.~\cite{zurItsOwlNumbers2025} show that finetuning on teacher-generated data transfers preferences.
    It would be interesting to check whether the same can already be achieved by simply finetuning on a corpus with elevated token frequency, without any teacher that already holds the preference.

    \subsubsection{Targeted Preference Injection}
    Our results show that the entanglement structure is rich enough to enable fine-grained control over model preferences.
    We observe that some tokens simultaneously increase the probability of an animal being mentioned as the most loved and the most hated, which at first appears contradictory.
    This is possible because effects on positive and negative emotions are uncorrelated, meaning a token's influence on most-loved is independent of its influence on most-hated.
    A single token can therefore induce a complex, context-dependent shift in how a concept is perceived, changing multiple emotional relations at once rather than changing only one.

    Conversely, because the effects are independent, other tokens affect only one emotional relation.
    A token that increases the probability of an animal being most loved may leave its most-hated probability entirely unchanged, allowing interventions precise enough to shift a single emotional relation in isolation.
    Together, these two properties suggest that subliminal prompting may be able to inject a specific behavioral signature into the model, selecting which emotional relations to shift and which to leave unchanged.
    The next step would be to test whether a target behavioral signature can be achieved in practice.

\end{qualitymark}

\bibliographystyle{plain}
\bibliography{references}

\include{appendix}
\end{document}